\scanpage{56}
\mspage{22.}

\diagram{\begin{tikzcd}[ampersand replacement=\&]
Y \arrow[r,phantom,"{\underset{\unsure{i}}{\subset}}" description] \arrow[dr,"g"'] \&
X \arrow[d,"f"]\\
\& \unsure{S}
\end{tikzcd}}

\uline{Th. de pureté (\deleted{\unsure{absolu}}\inserted{relatif})} comme conséquence du
th de dualité globale. Supposons que
$g$ soit \deleted{quasi-fini et}
\inserted{un \uline{isom}},
\deleted{\unsure{lisse} et \unsure{pour} \unsure{f}}
\inserted{[donc $i$ une imm. \uline{fermée}]}
alors la formule de dualité \deleted{\unsure{suivante}}
\[
\mathbb{R}_!\unsure{f_*|g_*}
\left(\underbrace{\mathbb{R}\underline{\operatorname{Hom}}(A_Y,\mathbb{R}^!f(G^\bullet))}
_{\mathbb{R}i_*\mathbb{R}i^!(\mathbb{R}^!f(\unsure{G^\bullet}))}\right)
\simeq \mathbb{R}\underline{\operatorname{Hom}}(\mathbb{R}_!g(A_Y),\deleted{\unsure{C}}G^\bullet)
\]
s'écrit
\[
g_*(\mathbb{R}i^!(\mathbb{R}^!f(G^\bullet)))\xrightarrow{\sim}G^\bullet
\]
donc si $G^\bullet$ est \unsure{réduit} au degré $0$, on
trouve en particulier
\[
\left\{
\begin{aligned}
\underline{H}^i_Y(f^*(G))&=0 &&\text{si }i\ne 2d\\
\underline{H}^{2d}_Y(\unsure{f^*|f^!}(G))&\simeq g^*(G)\otimes T_{Y/X}^{-1}
\end{aligned}
\right.
\qquad\deleted{\unsure{si}}
\]

\deleted{Par \unsure{transitivité}, on en déduit}
Si maintenant $Y\subset X$ est \deleted{\unsure{et}} une imm.
telle que $Y,X$ soient tous les deux lisses
sur $S$, de dim \unsure{quelconques},
\inserted{et de codim. relative $d$ dans $X$,}
\unsure{alors} \deleted{\unsure{pour} \unsure{être}}
\unsure{dès qu'on a} les $\simeq$ \unsure{voulus} pour \unsure{ces} \unsure{immers.}
Comme on a évidemment une flèche \unsure{can.}
$\longleftarrow$ par \unsure{cup-produits}, on voit que la question est
locale, or localement sur $X$ on peut trouver
un diagramme cartésien
\diagram{\begin{tikzcd}[ampersand replacement=\&]
Y \arrow[rr] \arrow[d] \&\& X \arrow[d]\\
Y' \arrow[rr,no head] \arrow[dr,"g'"',"\unsure{\sim}"] \&\&
X'=S[t_1,\unsure{\ldots|--},t_d] \arrow[dl,"f'"]\\
\& S \&
\end{tikzcd}}
\deleted{et \unsure{avec} \unsure{f}}
avec $X\longrightarrow X'$ \uline{lisse},
\mspage{55}
